\section{Why probabilistic semantics?}
Pure derivability is not enough for realistic parsing. A sentence may have several syntactically
well-formed derivations, and those derivations may correspond to different readings. This is not a
pathology but a normal feature of natural language. A parser that returns every derivation with
equal status is therefore mathematically respectable but practically incomplete.

Probabilistic grammatical models address this by associating weights or probabilities with lexical
choices and structural combinations. In probabilistic context-free grammars, probabilities rank
parse trees generated by grammar rules; in probabilistic CCG, similar ideas are applied to a more
lexicalised grammar formalism \citep{ManningSchutze1999,ClarkCurran2007}. There has also been
explicit prior work on stochastic Lambek grammars \citep{BonfanteDeGroote2004}. The basic
intuition is simple: if several analyses are licensed, prefer the one better supported by the data.

This thesis adopts that intuition in a deliberately conservative way. Rather than starting from a
full statistical parser, it starts from a semantic valuation $v(x,A) \in [0,1]$ for strings $x$ and
Lambek types $A$, interpreted relative to a corpus. The main question is then whether the
algebraic behaviour of the Lambek connectives can be captured in a way that supports a
sound proof theory.

